% Part: incompleteness
% Chapter: theories-computability
% Section: introduction

\documentclass[../../../include/open-logic-section]{subfiles}

\begin{document}

\olfileid{inc}{tcp}{int}

\olsection{परिचय}

\begin{editorial}
हा विभाग पुन्हा लिहिणे आवश्यक आहे.
\end{editorial}

आपल्याकडे पुढील गोष्टी उपलब्ध आहेत:
\begin{enumerate}
\item एखादे फलन $\Th{Q}$ मध्ये निरूपणीय असणे म्हणजे काय,
  याची व्याख्या (\olref[req][int]{defn:representable-fn})
\item एखादा संबंध $\Th{Q}$ मध्ये निरूपणीय असणे म्हणजे काय,
  याची व्याख्या (\olref[req][rel]{defn:representing-relations})
\item $\Th{Q}$ मध्ये निरूपणीय फलने नेमकी संगणनक्षम
  फलनेच असतात, असे सांगणारे प्रमेय
  (\olref[req][int]{thm:representable-iff-comp})
\item $\Th{Q}$ मध्ये निरूपणीय संबंध नेमके संगणनक्षम
  संबंधच असतात, असे सांगणारे प्रमेय
  (\olref[req][rel]{thm:representing-rels})
\end{enumerate}

{ \em उपपत्ती} म्हणजे तार्किक निष्पन्नतेखाली बंद असलेला वाक्यसंच.
म्हणजे $T$ ने $!A$ सिद्ध केले, तर $!A$ हे $T$ मध्ये
असते. उपपत्तीचा विचार वाक्यांचा आणि त्यांच्यापासून
निष्पन्न होणाऱ्या सर्व वाक्यांचा संग्रह म्हणून करणे
उपयुक्त ठरेल. यापुढे $\Th{Q}$ हे चिन्ह
\olref[req][int]{sec} मधील आठ स्वयंसिद्धकांपासून
निष्पन्न होणाऱ्या सर्व वाक्यांच्या { \em उपपत्ती}साठी
वापरू. $\Th{Q}$ मधील सूत्रांना संख्यांनी
सांकेतिक करता येते, हे लक्षात ठेवा. $!A$ असे
सूत्र असेल, तर $!A$ ला सांकेतिक करणारी
संख्या $\Gn{!A}$ ने दर्शवू. या सांकेतीकरणाच्या
आधारे आता सूत्रांचे विविध संच संगणनक्षम आहेत की
नाहीत, हा प्रश्न विचारता येतो.

\end{document}
